% Options for packages loaded elsewhere
\PassOptionsToPackage{unicode}{hyperref}
\PassOptionsToPackage{hyphens}{url}
\PassOptionsToPackage{dvipsnames,svgnames,x11names}{xcolor}
\documentclass[
  10pt,
  fontset=fandol]{ctexart}
\usepackage{xcolor}
\usepackage[margin=16mm]{geometry}
\usepackage{amsmath,amssymb}
\setcounter{secnumdepth}{-\maxdimen} % remove section numbering
\usepackage{iftex}
\ifPDFTeX
  \usepackage[T1]{fontenc}
  \usepackage[utf8]{inputenc}
  \usepackage{textcomp} % provide euro and other symbols
\else % if luatex or xetex
  \usepackage{unicode-math} % this also loads fontspec
  \defaultfontfeatures{Scale=MatchLowercase}
  \defaultfontfeatures[\rmfamily]{Ligatures=TeX,Scale=1}
\fi
\usepackage{lmodern}
\ifPDFTeX\else
  % xetex/luatex font selection
\fi
% Use upquote if available, for straight quotes in verbatim environments
\IfFileExists{upquote.sty}{\usepackage{upquote}}{}
\IfFileExists{microtype.sty}{% use microtype if available
  \usepackage[]{microtype}
  \UseMicrotypeSet[protrusion]{basicmath} % disable protrusion for tt fonts
}{}
\makeatletter
\@ifundefined{KOMAClassName}{% if non-KOMA class
  \IfFileExists{parskip.sty}{%
    \usepackage{parskip}
  }{% else
    \setlength{\parindent}{0pt}
    \setlength{\parskip}{6pt plus 2pt minus 1pt}}
}{% if KOMA class
  \KOMAoptions{parskip=half}}
\makeatother
\setlength{\emergencystretch}{3em} % prevent overfull lines
\providecommand{\tightlist}{%
  \setlength{\itemsep}{0pt}\setlength{\parskip}{0pt}}
\usepackage{amsmath,amssymb,mathtools,needspace}
\xeCJKDeclareCharClass{CJK}{"2032,"0400->"04FF,"2160->"216B,"2460->"2473,"25A0,"25C7,"25CB,"25CF}
\IfFontExistsTF{Arial Unicode MS}{\xeCJKsetup{AutoFallBack=true}\setCJKfallbackfamilyfont{\CJKrmdefault}{Arial Unicode MS}}{}
\setcounter{secnumdepth}{0}
\setlength{\emergencystretch}{2em}
\allowdisplaybreaks
\usepackage{bookmark}
\IfFileExists{xurl.sty}{\usepackage{xurl}}{} % add URL line breaks if available
\urlstyle{same}
\hypersetup{
  pdftitle={习题},
  colorlinks=true,
  linkcolor={Maroon},
  filecolor={Maroon},
  citecolor={Blue},
  urlcolor={Blue},
  pdfcreator={LaTeX via pandoc}}

\title{习题}
\author{}
\date{}

\begin{document}
\maketitle

\subsection{习题}\label{ux4e60ux9898}

\begin{enumerate}
\def\labelenumi{\arabic{enumi}.}
\item
  就初值问题 \(y'=ax+b,y(0)=0\) 分别导出 Euler 方法和改进的 Euler 方法的近似解的表达式，并与准确解 \(y=\dfrac12ax^2+bx\) 相比较。
\item
  用改进的 Euler 方法解初值问题 \(\begin{cases}y'=x+y,\quad0<x<1,\\y(0)=1,\end{cases}\) 取步长 \(h=0.1\) 计算，并与准确解 \(y=-x-1+2\mathrm e^x\) 相比较。
\item
  用改进的 Euler 方法解 \(\begin{cases}y'=x^2+x-y,\\y(0)=0,\end{cases}\) 取步长 \(h=0.1\) 计算 \(y(0.5)\)，并与准确解 \(y=-\mathrm e^{-x}+x^2-x+1\) 相比较。
\item
  用梯形方法解初值问题 \(\begin{cases}y'+y=0,\\y(0)=1,\end{cases}\) 证明其近似解为 \(y_n=\left(\dfrac{2-h}{2+h}\right)^n\)，并证明当 \(h\to0\) 时，它收敛于原初值问题的准确解 \(y=\mathrm e^{-x}\)。
\item
  利用 Euler 方法计算积分 \(\displaystyle\int_0^x\mathrm e^{t^2}\,\mathrm dt\) 在点 \(x=0.5,1,1.5,2\) 的近似值。
\end{enumerate}

\begin{quote}
校注（agent 补充，ch05b-E007）：小结中的五处节号按原印刷保留为 7.3.1、7.3.4、7.2.1、7.2.4、7.5.7；这些节号与本版章节编排不符。对应主题在本书章节索引中分别为 5.3.1（Taylor 级数法）、5.3.5（四阶 Runge-Kutta 方法）、5.2.1（Euler 格式）、5.2.4（改进的 Euler 格式）、5.5.7（Hamming 格式）。其中本章 PDF 140 的正文也明确将四阶 Runge-Kutta 格式指向 5.3.5，PDF 144 的实际标题为 5.5.7 Hamming 格式。
\end{quote}

\href{../../source-images/pdf-154.jpeg}{原页}

\textbf{6.} 取 \(h=0.2\)，用经典的四阶 Runge-Kutta 方法求解下列初值问题：

（1）\(\begin{cases}y'=x+y,\quad0<x<1,\\y(0)=1;\end{cases}\)

（2）\(\begin{cases}y'=3y/(1+x),\quad0<x<1,\\y(0)=1.\end{cases}\)

\textbf{7.} 证明对任意参数 \(t\)，下列 Runge-Kutta 格式是二阶的：

\[
\begin{cases}
\displaystyle y_{n+1}=y_n+\frac h2(K_2+K_3),\\[6pt]
K_1=f(x_n,y_n),\\
K_2=f(x_n+th,y_n+thK_1),\\
K_3=f(x_n+(1-t)h,y_n+(1-t)hK_1).
\end{cases}
\]

\textbf{8.} 证明下列两种 Runge-Kutta 方法是三阶的：

（1）

\[
\begin{cases}
\displaystyle y_{n+1}=y_n+\frac h4(K_1+3K_3),\\[6pt]
K_1=f(x_n,y_n),\\
\displaystyle K_2=f\left(x_n+\frac h3,y_n+\frac h3K_1\right),\\[4pt]
\displaystyle K_3=f\left(x_n+\frac23h,y_n+\frac23hK_2\right);
\end{cases}
\]

（2）

\[
\begin{cases}
\displaystyle y_{n+1}=y_n+\frac h9(2K_1+3K_2+4K_3),\\[6pt]
K_1=f(x_n,y_n),\\
\displaystyle K_2=f\left(x_n+\frac h2,y_n+\frac h2K_1\right),\\[4pt]
\displaystyle K_3=f\left(x_n+\frac34h,y_n+\frac34hK_2\right).
\end{cases}
\]

\textbf{9.} 分别用二阶显式 Adams 方法和二阶隐式 Adams 方法解下列初值问题：

\[
y'=1-y,\quad y(0)=0,
\]

取 \(h=0.2,y_0=0,y_1=0.181\)，计算 \(y(1.0)\) 并与准确解 \(y=1-\mathrm e^{-x}\) 相比较。

\textbf{10.} 证明下列解 \(y'=f(x,y)\) 的差分公式

\[
y_{n+1}=\frac12(y_n+y_{n-1})+\frac h4(4y'_{n+1}-y'_n+3y'_{n-1})
\]

是二阶的，并求出截断误差的首项。

\textbf{11.} 导出具有下列形式的三阶方法：

\[
y_{n+1}=a_0y_n+a_1y_{n-1}+a_2y_{n-2}+h(b_0y'_n+b_1y'_{n-1}+b_2y'_{n-2}).
\]

\textbf{12.} 将下列方程化为一阶方程组：

（1）\(y''-3y'+2y=0,\quad y(0)=1,\quad y'(0)=1\)；

（2）\(y''-0.1(1-y^2)y'+y=0,\quad y(0)=1,\quad y'(0)=0\)；

（3）\(x''(t)=-\dfrac{x}{r^3},\allowbreak \quad y''(t)=-\dfrac{y}{r^3},\allowbreak \quad r=\sqrt{x^2+y^2},\allowbreak \quad x(0)=0.4,\allowbreak \quad x'(0)=0,\allowbreak \quad y(0)=0,\allowbreak \quad y'(0)=2\)。

\textbf{13.} 取 \(h=0.25\)，用差分法解边值问题 \(\begin{cases}y''+y=0,\\y(0)=0,\quad y(1)=1.68.\end{cases}\)

\textbf{14.} 对方程 \(y''=f(x,y)\) 可建立差分公式 \(y_{n+1}=2y_n-y_{n-1}+h^2f(x_n,y_n)\)，试用这一公式求解初值问题 \(\begin{cases}y''=1,\\y(0)=y(1)=0,\end{cases}\) 验证计算解恒等于准确解 \(y(x)=\dfrac{x^2-x}{2}\)。

\textbf{15.} 取 \(h=0.2\)，用差分方法解边值问题 \(\begin{cases}(1+x^2)y''-xy'-3y=6x-3,\\y(0)-y'(0)=1,\quad y(1)=2.\end{cases}\)

\begin{quote}
校注（agent 补充，ch05b-E008）：习题 14 原书写``初值问题''，已照录；但给定条件为两端的 \(y(0)=y(1)=0\)，按本章定义属于边值问题，疑为称谓误植。此处只标记原书疑误，未代做习题。
\end{quote}

\subsection{本节覆盖原页的校注记录}\label{ux672cux8282ux8986ux76d6ux539fux9875ux7684ux6821ux6ce8ux8bb0ux5f55}

下列记录按原页关联，含同页相邻小节的校注；计算时同时读取上文校注和完整原页。

\begin{itemize}
\tightlist
\item
  \texttt{ch05b-E007}：原书 PDF 页 153
\item
  \texttt{ch05b-E008}：原书 PDF 页 154
\end{itemize}

\end{document}
